\documentclass[12pt,a4paper]{article}

\usepackage[a4paper,margin=18mm]{geometry}
\usepackage[T2A]{fontenc}
\usepackage[utf8]{inputenc}
\usepackage[russian]{babel}
\usepackage{amsmath,amssymb,mathtools}
\usepackage{array,booktabs,longtable}
\usepackage{xcolor}
\usepackage{hyperref}
\usepackage{tikz}

\hypersetup{
  colorlinks=true,
  linkcolor=blue!55!black,
  urlcolor=blue!55!black
}

\setlength{\parindent}{0pt}
\setlength{\parskip}{0.55em}
\renewcommand{\arraystretch}{1.25}

\begin{document}

\begin{titlepage}
\thispagestyle{empty}
\begin{center}
\vspace*{0.23\textheight}

{\LARGE\bfseries Ревью домашней работы\par}

\vspace{2.4cm}

{\large Автор: Лёвин Артём Александрович\par}

\vspace{0.8cm}

{\large 9 сентября 2026 года\par}

\vfill

\begin{tikzpicture}[x=1cm,y=1cm]
  \draw[blue!30,line width=0.8pt]
    (0,0) .. controls (2.2,0.45) and (4.4,-0.45) .. (6.6,0)
          .. controls (8.8,0.45) and (11.0,-0.45) .. (13.2,0);
\end{tikzpicture}

\vspace{1.0cm}
\end{center}
\end{titlepage}

\newpage
\section*{Сводная таблица}

{\small
\begin{longtable}{@{}>{\centering\arraybackslash}p{0.08\linewidth}p{0.19\linewidth}p{0.31\linewidth}p{0.14\linewidth}p{0.16\linewidth}@{}}
\toprule
\textbf{№ задания} & \textbf{Тема} & \textbf{Суть ошибки} & \textbf{Ваш ответ} & \textbf{Правильный ответ} \\
\midrule
\endfirsthead
\toprule
\textbf{№ задания} & \textbf{Тема} & \textbf{Суть ошибки} & \textbf{Ваш ответ} & \textbf{Правильный ответ} \\
\midrule
\endhead
\bottomrule
\endfoot
\hyperref[task:8]{8} &
Показательные уравнения: приведение к одному основанию &
После верной записи $\sqrt{125}=125^{1/2}$ не использовано равенство $125=5^3$. Из-за этого уравнение не доведено до равенства степеней с основанием $5$ и сделан неверный вывод об отсутствии решения. &
«нету» &
$x=1$ \\
\end{longtable}
}

\newpage
\vspace*{\fill}
\begin{center}
{\LARGE\bfseries Разбор ошибок}
\end{center}
\vspace*{\fill}
\newpage

\phantomsection\label{task:8}
\section*{Задание № 8}

\subsection*{Текст задания}
Решить уравнение
\[
\frac{5^{x+2}}{\sqrt{125}}=5^{\frac32}.
\]

\subsection*{Ваше решение}
В работе записан переход
\[
\frac{5^{x+2}}{\sqrt{125}}=5^{\frac32}
\quad\Longrightarrow\quad
\frac{5^{x+2}}{125^{\frac12}}=5^{\frac32}.
\]
После этого дальнейшее преобразование не выполнено; рядом с последней строкой поставлен знак вопроса.

\subsection*{Ваш ответ}
«нету».

\subsection*{Ошибка}
\textcolor{red}{Ошибка:} неверный вывод об отсутствии решения после незавершённого преобразования знаменателя.

\subsection*{Где именно возникает ошибка}
Последний корректный шаг:
\[
\sqrt{125}=125^{\frac12}.
\]
Это равенство верно: арифметический квадратный корень из положительного числа можно записать как степень с показателем $\frac12$.

Первый неверный по смыслу шаг возникает дальше: преобразование останавливается, а затем записывается ответ «нету». Само уравнение к этому моменту ещё не исследовано до конца, поэтому оснований утверждать, что решений нет, не было.

\subsection*{Почему это неверно}
Число $125$ связано с основанием $5$:
\[
125=5^3.
\]
Следовательно,
\[
125^{\frac12}=(5^3)^{\frac12}=5^{\frac32}.
\]
Именно это преобразование является ключевым: после него и числитель, и знаменатель, и правая часть записываются через одну и ту же основу $5$.

Для положительного основания $a\ne1$ функция $a^t$ принимает одинаковые значения только при одинаковых показателях. Поэтому, если получено равенство
\[
5^A=5^B,
\]
то можно приравнять показатели:
\[
A=B.
\]

В данном задании вывод «решений нет» сделан до применения этого правила. Наоборот, после приведения к одному основанию получается обычное линейное уравнение относительно показателя степени.

\subsection*{Ключевая идея}
\textcolor{blue!60!black}{Ключевая идея:} сначала представить постоянный знаменатель как степень того же основания, что и остальные степени:
\[
\sqrt{125}=\sqrt{5^3}=5^{\frac32}.
\]
После этого применить правило деления степеней с одинаковым основанием и приравнять показатели.

\subsection*{Исправление от последнего правильного шага}
Продолжим именно с верной записи
\[
\frac{5^{x+2}}{125^{\frac12}}=5^{\frac32}.
\]
Так как $125=5^3$, то
\[
125^{\frac12}=(5^3)^{\frac12}=5^{\frac32}.
\]
Поэтому
\[
\frac{5^{x+2}}{5^{\frac32}}=5^{\frac32}.
\]
При делении степеней с одинаковым основанием показатели вычитаются:
\[
5^{x+2-\frac32}=5^{\frac32}.
\]
Слева
\[
x+2-\frac32=x+\frac12,
\]
значит,
\[
5^{x+\frac12}=5^{\frac32}.
\]
Основания одинаковы и равны $5$, поэтому показатели равны:
\[
x+\frac12=\frac32.
\]
Отсюда
\[
x=\frac32-\frac12=1.
\]

\subsection*{Полное правильное решение}
Рассмотрим уравнение
\[
\frac{5^{x+2}}{\sqrt{125}}=5^{\frac32}.
\]

Сначала преобразуем знаменатель. Поскольку
\[
125=5^3,
\]
имеем
\[
\sqrt{125}=\sqrt{5^3}=(5^3)^{\frac12}=5^{\frac32}.
\]
Знаменатель положителен и не равен нулю, поэтому ограничений на $x$ из знаменателя нет.

Подставляем полученное выражение:
\[
\frac{5^{x+2}}{5^{\frac32}}=5^{\frac32}.
\]
Используем правило
\[
\frac{a^m}{a^n}=a^{m-n},\qquad a\ne0:
\]
\[
5^{x+2-\frac32}=5^{\frac32}.
\]
Упростим показатель слева:
\[
2-\frac32=\frac12,
\]
поэтому
\[
5^{x+\frac12}=5^{\frac32}.
\]
Так как основание $5$ положительно и $5\ne1$, приравниваем показатели:
\[
x+\frac12=\frac32.
\]
Вычитаем $\frac12$ из обеих частей:
\[
x=1.
\]

\subsection*{Проверка}
Подставим $x=1$ в левую часть исходного уравнения:
\[
\frac{5^{1+2}}{\sqrt{125}}
=\frac{5^3}{5^{\frac32}}
=5^{3-\frac32}
=5^{\frac32}.
\]
Получено ровно то же выражение, что стоит в правой части. Следовательно, $x=1$ действительно является решением.

\subsection*{Итог}
\textcolor{teal!60!black}{Итог:} решение уравнения существует и единственно:
\[
\boxed{x=1}.
\]
Ошибка состояла не в переходе $\sqrt{125}=125^{1/2}$, а в том, что преобразование не было продолжено до $125=5^3$, после чего был сделан необоснованный вывод об отсутствии решения.

\subsection*{Задача на закрепление}
Решите уравнение
\[
\frac{3^{x+1}}{\sqrt{27}}=3^{\frac52}.
\]

\end{document}
